% Canonical node 011 · C013; current section path 1.2.4.
% Canonical owner: C013
\cnnaNodeHeading{011}{C013}{First non-root provenance birth $v_1$}
\cnnaNodePlacement{1.2.4}{D1}

\paragraph*{Position and scientific role.}
\texttt{C013} joins three independent bootstrap inputs: the structural first
slot \texttt{C004A}, the singleton token \texttt{A001}, and the fixed unit
normalization \texttt{N001}.  It is exceptional because before this step there
is no nontrivial weighted relation and hence no Schur/DtN response from which a
response-coupled birth could be computed.

% CNNA-ARCHITECTURE-BEGIN C013
\begin{cnnaInfoBox}{First nontrivial ComplemeNt Net}
C013 realizes the first selected open slot and joins the newborn to the existing root by the first weighted directed relation.  This is the first nontrivial \emph{net} in CNNA.  Complementarity is relational and stage-dependent: root and newborn are not declared permanent opposites, but jointly form the first realized response context from which later relative cuts can be taken.
\end{cnnaInfoBox}
% CNNA-ARCHITECTURE-END C013

\paragraph*{Mathematical contract and construction.}
Let
\begin{equation}
  s_1=(\varepsilon,0,(0))
\end{equation}
be the C004A slot and assume its address lies in the finite cutoff.  C013
constructs
\begin{equation}
  v_1=(0),\qquad \operatorname{parent}(v_1)=r=\varepsilon,
\end{equation}
with the two stored directed relations and conductances
\begin{equation}
  \mathcal R_1=
  \bigl((r,v_1),(v_1,r)\bigr),
  \qquad
  \mathcal C_1=(1,1).
  \label{eq:c013-first-birth}
\end{equation}
The output stores the slot, the normalization and, in Lean, the proof of cutoff
admission.  The A001 token is consumed as a constructor argument but is not a
field of the resulting state.

\paragraph*{Invariants and boundary.}
The newborn address is definitionally the C004A address, its syntactic parent
is the root, and both conductances are inherited from N001.  No event index,
birth time, geometry, response, steering value or later live-state update is
created.  Since $|(0)|=1$, the construction exists for $L\ge1$ and is rejected
at $L=0$.  The structural word $(0)$ continues to exist at $L=0$; only its
promotion to a born carrier element fails.

\paragraph*{Cross-layer realization and counterchecks.}
Python checks exact predecessor types and executes the cutoff guard before
constructing an immutable two-field birth object.  Lean moves the cutoff guard
into the input proposition and stores its proof in the structure.  Both layers
derive the endpoints, relation orientations and unit pair from the same
predecessors.  The Python tests check the complete first relation, the $L=0$
rejection and the absence of a retained seed field.  These checks exclude both
an illicit response-driven bootstrap and hidden information transfer from A001.

\paragraph*{Canonicity boundary, result and handoff.}
For the fixed C004A slot and N001 normalization, the endpoint and conductance
data are canonical.  C013 does not itself state the equality of outputs under
two explicit seed arguments; \texttt{T001} owns that theorem.  Verified edge
\texttt{E009} sends the construction to T001 and verified edge \texttt{E011}
sends the completed birth to \texttt{C014}.

\noindent\textbf{Primary code anchors.}
Python construction: lines 23--63; focused tests: lines 17--33.  Lean structure,
constructor, derived endpoints, relation pair, conductance pair and three
invariant theorems: lines 22--77.  The Supplement records all 13 anchors and
source hashes.
